\documentclass[10pt,a4paper]{article}
\usepackage{fontspec,amsmath,amssymb,amsthm,mathtools,graphicx,bm}
\usepackage[margin=22mm]{geometry}
\usepackage[bookmarks=false]{hyperref}
\setlength{\emergencystretch}{3em}
\title{Width-preserving lifting from forgetful parity decision trees for strongly stifled gadgets}
\author{Arkadev Chattopadhyay, Pavel Dvořák}
\date{doi:10.4230/LIPIcs.CCC.2025.24; CC BY 4.0}
\begin{document}
\maketitle
\noindent\textit{Editorial note.} Editable TeX faithfully converted from the official CC BY LaTeXML HTML prose and math alttext; the official PDF/HTML originals are included. Only the main Theorem 13 lifting/simulation result counts. Its strongly stifled gadget hypothesis, canonical tree model, query-depth bound and width bound are retained; no unrestricted proof-system lower bound is claimed.\par\medskip
\noindent\textit{Allowed standard prerequisites (editorial).}\begin{itemize}
\item Standard propositional resolution and finite-field linear algebra. Source-local definitions and all author support proofs retained.
\end{itemize}\medskip


\par\medskip\noindent\textbf{
1 Introduction}\par




Understanding trade-offs among complexity measures in a computational model is a well known interesting theme, with many published results (for example, time-space trade-offs [15, 16, 26], rounds-communication trade-offs [21, 10, 2] and space-size trade-offs in propositional proof complexity [6, 3, 5]). Typically, in these trade-offs, one showed that in various models of computation, simultaneous optimization of two complexity measures, like space and time, or rounds and total communication, or space and width (in refuting CNF formulas) is not always possible. In particular trying to optimize one complexity measure, necessarily leads to a huge blow-up in the other measure. For instance, in Yao’s 2-party model of communication, the Greater-Than function can be computed in $1$ round. It can also be computed using randomized protocols of communication cost $O(\log n)$. But every $O(1)$-round protocol, requires $\Omega(n)$ communication cost. On the other hand, every function has a protocol of cost $O(n)$. In all of the trade-off results cited above, the general story was that trying to optimize the use of one resource, led to the cost with respect to to the other resource shooting up to the cost needed by a naive/generic algorithm.\par




In 2016, Razborov [24] exhibited formulas for which very different and extreme kind of trade-offs hold in the propositional proof system of resolution. Although these unsatisfiable $k$-CNF formulas on $n$ variables have refutations of $O(k)$ width, every one of their tree-like refutation of width less than $n^{1-\epsilon}/k$ has size $\text{exp}(n^{\Omega(k)})$. That is, despite the fact that every $n$-variable formula has a generic tree-like refutation of size $2^{n}$, these exhibited formulas that do have refutation of small width require super-critical tree-like refutation size whenever width is mildly restricted. Moreover, the super-critical size is in fact exponentially larger than the generic upper bound. Razborov remarked that such a phenomenon seemed extremely rare in the known body of tradeoff results in the computational complexity literature. In concluding his work, he urged finding more instances of such trade-offs. In response to that, follow-up works have appeared. They can be classified into two types. Ones which continue to focus on resolution and, others on more powerful proof systems. Examples of the former include work by Berkholz and Nordstrom [8], who showed super-critical trade-offs between width and space. A recent work of Berkholz, Lichter and Vinall-Smeeth [7] proves super-critical trade-offs for narrow resolution width and narrow tree-like size for refuting the isomorphism of two graphs.\par




The second type of work answers Razborov’s call by finding such trade-offs in stronger proof systems. This includes the recent work of Fleming, Pitassi and Robere [14] who first showed that the argument of Razborov extends to general resolution DAGs. They then use it along with appropriate lifting theorems to prove trade-offs between size and depth for DAG like Resolution, Res$(k)$, and cutting planes.
In very recent progress in the area,
De Rezende, Fleming, Jannet, Nordström, and Pang [12], and Göös, Maystre, Risse, Sokolov [18] independently showed super-critical trade-offs not only for various proof systems but also the first super-critical depth-size tradeoffs for monotone circuits.
Our work also falls in this second type as we study tree-like resolution over parities, that generalizes tree-like resolution.\par




We exhibit super-critical trade-offs for width and tree-like size/depth in the style of Razborov for resolution over parities, denoted by $\text{Res}(\oplus)$. This system, introduced by Itsykson and Sokolov [19, 20], is one of the simplest generalizations of resolution for which obtaining super-polynomial lower bounds on size of refutations is a current, well known, challenge. Very recent works (see [13, 9]) managed to obtain exponential lower bounds on the size of \emph{regular proofs} in this system.\par




Our work here will concern tree-like $\text{Res}(\oplus)$ proofs. Lower bounds for them were obtained by Itsykson and Sokolov [20] themselves. More recently, two independent works, one by Beame and Koroth [4] and the other by Chattopadhyay, Mande, Sanyal and Sherif [11], proved lifting theorems that yielded a systematic way of lifting tree-like resolution width complexity to strong lower bounds on size of tree-like $\text{Res}(\oplus)$ proofs for formulas lifted with constant-size gadgets.
In this paper, we extend the lifting theorem by Chattopadhyay et al. [11] in the following manner. Their result was applicable to parity decision trees (duals of tree-like $\text{Res}(\oplus)$ proofs) that only had usual nodes where the algorithm queried (correspondingly the proof resolved on) an $\mathbb{F}_{2}$ linear form. We call such nodes query nodes. On the other hand, we want to deal here with width-bounded proofs that could be much deeper than $n$, the total number of variables of the formula. This would correspond to parity decision trees where the height is much larger than $n$, and therefore, necessarily there are nodes that \emph{forget}. The affine space corresponding to such a forget node $u$ is strictly contained in the affine space corresponding to $u$’s only child node $v$. Alternatively, in the bottom-up view of the corresponding proof, the linear clause at $v$ is strictly weakened to get the linear clause at $u$. Dealing with such nodes, so that the width of the (ordinary) clauses in the extracted resolution proof never exceed the corresponding width of the linear clauses, is the main technical contribution of this work.
Thus, we establish a depth-to-size lifting result from tree-like $\text{Res}(\oplus)$ of arbitrary depth to tree-like resolution, which also preserves the width of the refutation.\par



% SOURCE-ID Thmtheorem1


\par\medskip\noindent\textbf{
\textbf{Theorem 1}\textbf{.}
}\par



\emph{Let $\Phi\circ g$ be a lift of a contradiction $\Phi$ by an appropriate gadget $g:\{0,1\}^{\ell}\to\{0,1\}$.
Suppose there is a tree-like $\text{Res}(\oplus)$ refutation for $\Phi\circ g$ with size $s$ and width $w$.
Then, there is a tree-like resolution refutation for $\Phi$ with depth at most $\log s$ and width at most $w$.}\par



% END-SOURCE-ID Thmtheorem1


% SOURCE-ID Thmtheorem2


\par\medskip\noindent\textbf{
$\blacktriangleright$ Remark 2.
}\par



\emph{We point out the precise difference between our Theorem 1 and the earlier lifitng theorem of Chattopadhyay et al [11]. The earlier theorem, given a tree-like refutation of $\Phi\circ g$ in $\text{Res}(\oplus)$ of size $s$ and width $w$, would have extracted a tree-like refutation of $\Phi$ in ordinary resolution of depth $\log s$, with no guarantees on the width of this refutation. In fact, the width could get as large as the depth of the extracted refutation, i.e. $\log s$. In super-critical trade-offs, which is our chief interest here, the width $w$ of the given $\text{Res}(\oplus)$ refutation of $\Phi\circ g$ could be exponentially smaller than $\log s$. This renders the earlier lifting theorem unusable for demonstrating such trade-offs.}\par



% END-SOURCE-ID Thmtheorem2



Applying Theorem 1 to the trade-off by Razborov [24], we immediately obtain an analogous trade-off in the $\text{Res}(\oplus)$ proof system.\par



% SOURCE-ID Thmtheorem3


\par\medskip\noindent\textbf{
\textbf{Theorem 3}\textbf{.}
}\par



\emph{Let $k=k(n)\geq 12$ be any parameter and let $\varepsilon>0$ be an arbitrary constant.
Then, there exists a $k$-CNF contradiction $\tau$ over $n$ variables such that there is a resolution refutation for $\tau$ with width at most $O(k)$, but for every tree-like $\text{Res}(\oplus)$ refutation $\Pi$ for $\tau$ with $w(\Pi)\leq n^{1-\varepsilon}/k$, we have the bound $|\Pi|\geq\exp(n^{\Omega(k)})$.}\par



% END-SOURCE-ID Thmtheorem3



The contradiction $\tau$ from the previous theorem is a lift of the contradiction $\tau^{\prime}$ constructed by Razborov [24] by an appropriate gadget $g:\{0,1\}^{\ell}\to\{0,1\}$ with a constant size.
A caveat of $\tau^{\prime}$ (as Razborov also noted) is that the number of clauses of $\tau^{\prime}$ is $n^{\Theta(k)}$.
Naturally, this caveat is inherited by our contradiction $\tau$.
This issue was addressed in very recent work of de Rezende et al. [12].
They provided a contradiction $\chi$ such that the size of its tree-like resolution with bounded width is super-exponential in not just in the number of variables but also in the size of the formula.
However, the bound on width of the tree-like resolution for which super-critical size is needed is much more strict than in Razborov’s result.
They showed the super-critical trade-off only for tree-like resolution of width smaller than $2w$, where $w$ is the width of the resolution refutation of $\chi$.
Since our gadget has size 3, we can guarantee only resolution refutation with width at most $3w$ for the lifted formula $\chi\circ g$.
Thus, we can not lift their super-critical trade-off to $\text{Res}(\oplus)$ as it is extremely sensitive to width.
Razborov’s result [24], on the other hand, is more robust making it possible to be lifted by our Theorem 1.\par




If one were able to construct another formula $\tilde{\tau}$ improving state-of-the-art in supercritical trade-off between width and size of tree-like resolution that is not so sensitive to the width of the refutations, this improvement could be used with our simulation theorem (Theorem 1) and be lifted to tree-like $\text{Res}(\oplus)$.\par




\par\medskip\noindent\textbf{Relation to Other Recent Works}\par




The $\text{Res}(\oplus)$ proof system has been an active area of research.
Recently, Efremenko, Garlík, and Itsykson [13] showed that the binary pigeonhole principle formula requires an exponential-size refutation within the so-called bottom-regular $\text{Res}(\oplus)$.
The bottom-regular $\text{Res}(\oplus)$ is a fragment of $\text{Res}(\oplus)$ that contains both tree-like $\text{Res}(\oplus)$ and regular resolution proof systems.
Furthermore, Bhattacharya, Chattopadhyay, and Dvořák [9] showed that bottom-regular $\text{Res}(\oplus)$ can not polynomially simulate even ordinary, but DAG-like resolution.
This separation was very recently improved quantitatively by Alekseev and Itsykson [1].\par




Furthermore, Alekseev and Itsykson [1] established a width-to-width lifting from resolution to $\text{Res}(\oplus)$.
They proved this in a contra-positive way – if there is no resolution refutation of a contradiction $\Phi$ with width $w$, then there is no width-$w$ $\text{Res}(\oplus)$ refutation of a lift of $\Phi$ by an appropriate gadget $g$.
They utilized a game interpretation of resolution and $\text{Res}(\oplus)$ to prove their lifting theorem.
While their proof is quite short, it is unclear whether their technique can be adapted to prove the depth-to-size lifting theorem as we need in order to show the trade-off in $\text{Res}(\oplus)$ (our Theorem 3). In particular, their theorem seems incomparable to the depth-to-size lifting of Chattopadhyay et al. [11].
On the other hand, since any refutation can be expanded into a tree-like refutation (with a possible exponential blow-up in size), our lifting theorem (Theorem 1) immediately implies the width-to-width lifting theorem of Alekseev and Itsykson (however our proof seems more involved). Hence, our Theorem 1 effectively contains a common generalization of the width-to-width lifting of Alekseev and Itsykson [1] and depth-to-size lifting of Chattopadhyay et al. [11].
Moreover, we use a completely different technique than Alekseev and Itsykson [1].
Specifically, we establish our lifting theorem directly by constructing a tree-like resolution refutation for a contradiction $\Phi$ simulating a tree-like $\text{Res}(\oplus)$ refutation for $\Phi\circ g$.
To achieve this, we use some ideas from linear algebra that, to our knowledge, have not been previously utilized in the context of lifting.\par





\par\medskip\noindent\textbf{Overview of Our Ideas}\par




Overall, the ideas behind the proof of Theorem 1 are inspired by the work of Chattopadhyay et al. [11].
However, they did not consider $\text{Res}(\oplus)$ refutation with a limited width.
Thus, they only needed to process query nodes to prove their lifting theorem.
In contrast, our setting also involves forget nodes, where a linear equation from the span of previously queried equations is forgotten.
It turns out that processing forget nodes is non-trivial. In particular, an affine space can be viewed in two ways: the first is by the (linear) space of constraints that could be thought of as the dual view. The primal view is that of the set of vectors lying in space represented by a basis of the underlying vector space and a shift vector. Previously, in [11], the dual view was very effectively used for depth-to-size lifting in the absence of forget nodes. This is because a query node naturally adds a constraint to the dual space. On the other hand, a forget node increases the dimension of the affine space. This new space is not conveniently representable with respect to the basis maintained for the dual space of constraints just before “forgetting” happens. Here it seems the primal view is more helpful as any basis of a space $A_{1}$ can be extended to a basis of a space $A_{2}$ whenever $A_{1}\subseteq A_{2}$. The main tool we use is a a characterization, via Theorem 10, of the constraint space of $A_{2}$ in terms of the constraint space of $A_{1}$, where $A_{1}\subseteq A_{2}\subseteq\mathbb{F}_{2}^{n}$. The proof of this turns out to be simple\footnote{Our original proof was complicated. The much simpler proof we present here has been pointed out to us by an anonymous referee.}. With more ideas, including a new notion of \emph{strongly stifled} gadgets that extends the earlier notion of stifling introduced by [11], Section 7 yields the process of dealing with forget nodes.\par






\par\medskip\noindent\textbf{
2 Tree-like Proofs and Decision Trees}\par




A proof in a propositional proof system starts from a set of clauses $\Phi$, called axioms, that is purportedly unsatisfiable. It generates a proof by deriving the empty clause from the axioms, using inference rules. The main inference rule in the standard resolution, called the resolution rule, derives a clause $A\vee B$ from clauses $A\vee x$ and $B\vee\lnot x$ (i.e., we resolve the variable $x$).
If we can derive the empty clause from the original set $\Phi$, then it proves that the set $\Phi$ is unsatisfiable.\par




Resolution over parities ($\text{Res}(\oplus)$) is a generalization of the standard resolution, using linear clauses (disjunction of linear equations in $\mathbb{F}_{2}$) to express lines of a proof.
It consists of two rules:\par


\textbf{Resolution Rule:}



From linear clauses $A\vee(\ell=0)$ and $B\vee(\ell=1)$ derive a linear clause $A\vee B$.\par



\textbf{Weakening Rule:}



From a linear clause $A$ derive a linear clause $B$ that is semantically implied by $A$ (i.e., any assignment satisfying $A$ also satisfies $B$).\par





The length $|\Pi|$ of a resolution (or $\text{Res}(\oplus)$) refutation $\Pi$ of a formula $\Phi$ is the number of applications of the rules above in order to refute the formula $\Phi.$
The width $w(\Pi)$ of a resolution (or $\text{Res}(\oplus)$) refutation $\Pi$ is the maximum width of any (linear) clause that is used in the resolution proof.
A (linear) resolution proof is \emph{tree-like} if the resolution rule is applied in a tree-like fashion.
The depth $d(\Pi)$ of the tree-like proof $\Pi$ is the depth of the underlying tree (i.e., the length of the longest path from the root to a leaf).\par




We can replace the general resolution rule with a canonical one:\par


\textbf{Canonical Resolution Rule:}



From linear clauses $C\vee(\ell=0)$ and $C\vee(\ell=1)$ derive a linear clause $C$.\par





Using the canonical resolution rule instead of the general one will not make the proof system substantially weaker.
If we want to apply the resolution rule on the clauses $A\vee(\ell=0)$ and $B\vee(\ell=1)$, we can apply the weakening rule to both of them to get linear clauses $A\vee B\vee(\ell=0)$ and $A\vee B\vee(\ell=1)$ and then apply the canonical resolution rule to derive the clause $A\vee B$.
Thus, from a tree-like $\text{Res}(\oplus)$ refutation $\Pi$ of a contradiction $\Phi$, we can derive an equivalent tree-like $\text{Res}(\oplus)$ refutation $\Pi^{\prime}$ that uses only canonical resolution rule (and the weakening rule) with $|\Pi^{\prime}|\leq 3\cdot|\Pi|$, $d(\Pi^{\prime})\leq 2\cdot d(\Pi)$, and $w(\Pi^{\prime})\leq w(\Pi)+1$ as for each application of the resolution rule in $\Pi$ we add 2 applications of the weakening rule in $\Pi^{\prime}$, and the width might increase by 1 as we introduce clause $A\vee B\vee(\ell=0)$ and $A\vee B\vee(\ell=1)$ but only the clause $A\vee B$ was present in $\Pi$.\par




It is known that a tree-like resolution (or $\text{Res}(\oplus)$) proof, for an unsatisfiable set of clauses $\Phi$, corresponds to a (parity) decision tree for a search problem $\textit{Search}(\Phi)$
that is defined as follows.
For a given assignment $\alpha$ of the $n$ variables of $\Phi$, one needs to find a clause in $\Phi$ that is not satisfied by $\alpha$ (at least one exists as the set $\Phi$ is unsatisfiable).
The correspondence holds even for general (not only tree-like) proofs
(see for example Garg et al. [17], who credit it to earlier work of Razborov [23] that was simplified by Pudlák [22] and Sokolov[25]), but in this paper, we are interested only in tree-like proofs.\par




Let $R\subseteq\{0,1\}^{m}\times O$, where $O$ is a set of possible outputs.
A forgetting parity decision tree (FPDT) computing $R$ is a tree $\mathcal{T}$ such that each node has at most two children and the following conditions hold:\par


\begin{itemize}


\item[] 
$\blacksquare$ 


Each node $v$ is associated with an affine space $A_{v}\subseteq\mathbb{F}_{2}^{m}$.\par





\item[] 
$\blacksquare$ 


For every node $v$ with two children $u$ and $w$ is called a \emph{query nodes}. There is a linear query $f_{v}$ such that $A_{u}=\{x\in A_{v}|\langle f_{v},x\rangle=0\}$ and $A_{w}=\{x\in A_{v}|\langle f_{v},x\rangle=1\}$, or vice versa.
We say that $f_{v}$ is the \emph{query} at $v$.\par





\item[] 
$\blacksquare$ 


Every node $v$ with exactly one child $u$ is called a \emph{forget node}. It holds that $A_{v}\subseteq A_{u}$.\par





\item[] 
$\blacksquare$ 


Each leaf $\ell$ is labeled by $o_{\ell}\in O$ such that for all $x\in A_{\ell}$, it holds that $(x,o_{\ell})\in R$.\par





\item[] 
$\blacksquare$ 


For the root $r$, $A_{r}=\mathbb{F}_{2}^{m}$.\par





\end{itemize}




The size $|\mathcal{T}|$ of an FPDT $\mathcal{T}$ is the number of nodes of $\mathcal{T}$ and the width $w(\mathcal{T})$ of FPDT $\mathcal{T}$ is the largest integer $w$ such that there exists an affine space of co-dimension at least $w$ associated with some node of $\mathcal{T}$.
The depth of $\mathcal{T}$ is denoted $d(\mathcal{T})$.
Note that there are no forget nodes in a standard parity decision tree.
Thus, for such trees, the width is exactly the depth of the tree.
It no longer holds for this model, because we may “forget” some linear queries that have been made earlier.\par




A forgetting decision tree (FDT) is defined similarly to FPDT but instead of affine spaces, cubes are associated with each node.
Consequently, the width $w(\mathcal{T})$ of FDT is the maximum number $w$ such that there exists a cube of width at least $w$ associated with some node of $\mathcal{T}$ and queries of single variables replace the linear queries at nodes.\par




The correspondence between an F(P)DT’s and tree-like resolution (or $\text{Res}(\oplus)$) proofs using only canonical resolution rule is the following:
We represent a (linear) resolution proof as a tree where nodes are associated with (linear) clauses.
The leaves are associated with clauses of $\Phi$ and the root is associated with the empty clause.
Each node with two children corresponds to an application of the canonical resolution rule and each node with exactly one child corresponds to an application of the weakening rule.
To get an F(P)DT for $\textit{Search}(\Phi)$ we just negate the clauses that are associated with the nodes.
Thus, each node is associated with a cube (or an affine space in the case of $\text{Res}(\oplus)$/FPDT).
Moreover, a cube $C_{v}$ (or an affine space $A_{v})$ associated with the node $v$ of an F(P)DT $\mathcal{T}$ contains exactly the falsifying assignments of the (linear) clause that is associated with the corresponding node in the tree-like refutation $\Pi$ that corresponds to $\mathcal{T}$.
It is clear that the width and the depth of such decision tree $\mathcal{T}$ are exactly the same as the width and the depth of the corresponding tree-like refutation $\Pi$ and the length $|\Pi|$ equals the number of inner nodes of $\mathcal{T}$ (as the inner nodes of $\mathcal{T}$ correspond to the applications of the resolution rule).\par




We say an FPDT $\mathcal{T}$ is \emph{canonical} if for each forget node $v$ of $\mathcal{T}$ and its only child $u$, it holds that $\text{co-dim}(A_{v})=\text{co-dim}(A_{u})+1$.
We say an FPDT $\mathcal{T}$ is \emph{succinct} if the parent of each forget node is a query node.
Note that any FPDT can be transformed into an equivalent canonical (or succinct) FPDT by expanding forget nodes into paths of forget nodes (or contracting paths of forget nodes to single vertices).\par




Consider an FPDT $\mathcal{T}$ and its succinct form $\bar{\mathcal{T}}$.
Note that the number of query nodes of $\bar{\mathcal{T}}$ and $\mathcal{T}$ is the same,
and analogously the number of query nodes on a root-leaf path in $\bar{\mathcal{T}}$ equals the number of query nodes of the corresponding path in $\mathcal{T}$.
Thus, for an FPDT $\mathcal{T}$ we define query size $|\mathcal{T}|_{\mathsf{q}}$ and query depth $d_{\mathsf{q}}(\mathcal{T})$ to be the number of query nodes of $\mathcal{T}$ and the maximum number of query nodes on a root-leaf path of $\mathcal{T}$.\par



% SOURCE-ID Thmtheorem4


\par\medskip\noindent\textbf{
\textbf{Observation 4}\textbf{.}
}\par



\emph{Let $\Pi$ be a $\text{Res}(\oplus)$ refutation of a contradiction $\Phi$.
Then, there is a canonical FPDT $\mathcal{T}$ computing $\textit{Search}(\Phi)$ with $|\mathcal{T}|_{\mathsf{q}}\leq|\Pi|$ and $w(\mathcal{T})\leq w(\Pi)+1$.}\par



% END-SOURCE-ID Thmtheorem4



\par\medskip\noindent\textbf{Proof.}\par




As discussed above, first we modify $\Pi$ to a $\text{Res}(\oplus)$ refutation $\Pi^{\prime}$ that uses only canonical resolution rule (and weakening rule).
By this modification, we have $w(\Pi^{\prime})\leq w(\Pi)+1$ and the number of applications of the resolution rule in $\Pi$ is exactly the number of applications of the canonical resolution rule in $\Pi^{\prime}$.
From $\Pi^{\prime}$, we derive an FPDT $\mathcal{T}^{\prime}$ computing $\textit{Search}(\Phi)$ that we finally modify to an equivalent canonical FPDT $\mathcal{T}$.
The width of $\Pi^{\prime}$ and $\mathcal{T}^{\prime}$ is the same, and the query size of $\mathcal{T}$ is exactly the number of applications of the canonical resolution rule in $\Pi^{\prime}$.
The modification of $\mathcal{T}^{\prime}$ to the canonical FPDT $\mathcal{T}$ does not change the width and the query size of the tree.
Thus, we have $w(\mathcal{T})\leq w(\Pi)+1$ and $|\mathcal{T}|_{\mathsf{q}}\leq|\Pi|$.
$\blacktriangleleft$\par



% SOURCE-ID Thmtheorem5


\par\medskip\noindent\textbf{
\textbf{Observation 5}\textbf{.}
}\par



\emph{Let $\mathcal{T}$ be a succinct FPDT computing $\textit{Search}(\Phi)$ and $\Pi$ be the tree-like $\text{Res}(\oplus)$ refutation of $\Phi$ corresponding to $\mathcal{T}$. Then, $|\Pi|\leq 3\cdot 2^{d_{\mathsf{q}}(\mathcal{T})}$.}\par



% END-SOURCE-ID Thmtheorem5




\par\medskip\noindent\textbf{Proof.}\par




As mentioned above, the length of $\Pi$ equals the number of inner nodes of $\mathcal{T}$.
The tree $\mathcal{T}$ has at most $2^{d_{\mathsf{q}}(\mathcal{T})}-1$ query nodes.
Since $\mathcal{T}$ is succinct, the number of forget nodes is at most twice the number of query nodes as each query node might have at most two forget nodes as its children. Further, any child of a forget node is not a forget node.
Thus, $|\Pi|\leq 3\cdot 2^{d_{\mathsf{q}}(\mathcal{T})}$.
$\blacktriangleleft$\par






\par\medskip\noindent\textbf{
3 Lifting of Relations and Formulas}\par




Let $g:\{0,1\}^{\ell}\to\{0,1\}$ be a boolean function.
For a relation $R\subseteq\{0,1\}^{n}\times O$ we define its lift $R\circ g\subseteq\{0,1\}^{n\ell}\times O$ as\par


\[\begin{gathered}
R\circ g=\bigl{\{}(y,o)\in\{0,1\}^{n\ell}\times O\mid\bigl{(}\overrightarrow{g%
}(y),o\bigr{)}\in R\bigr{\}},
\end{gathered}\]


where $\overrightarrow{g}(y^{1}_{1},\dots,y^{1}_{\ell},\dots,y^{n}_{1},\dots,y^{n}_{%
\ell})=\bigl{(}g(y^{1}_{1},\dots,y^{1}_{\ell}),\dots,g(y^{n}_{1},\dots,y^{n}_{%
\ell})\bigr{)}$.\par





For CNF $\Phi$ over $n$ variables $\{x_{1},\dots,x_{n}\}$, let $\Phi\circ g$ be the following lift of $\Phi$ over the variables $\{y^{i}_{j}\mid i\in[n],j\in[\ell]\}$.
For any clause $D$ of $\Phi$, let $\text{Vars}(D)$ be the set of variables of $D$, and let $\eta_{D}\in\{0,1\}^{\text{Vars}(D)}$ be the only falsifying assignment of $D$.
Then,\par


\[\begin{gathered}
D\circ g=\left\{\bigvee_{x_{i}\in\text{Vars}(D),j\in[\ell]}y^{i}_{j}\neq\kappa%
^{i}_{j}\nobreak\ \Bigl{|}\nobreak\ \kappa\in%
\overrightarrow{g}^{-1}(\eta_{D})\right\},
\end{gathered}\]


where $y^{i}_{j}\neq\kappa^{i}_{j}$ is the following literal:\par


\[\begin{gathered}
y^{i}_{j}\neq\kappa^{i}_{j}=\begin{cases}y^{i}_{j}&\text{ if }\kappa^{i}_{j}=0%
,\\
\neg y^{i}_{j}&\text{ if }\kappa^{i}_{j}=1.\end{cases}
\end{gathered}\]


Now, the clauses of $\Phi\circ g$ are $\{D\circ g\mid D\text{ clause of }\Phi\}$.\par



% SOURCE-ID Thmtheorem6


\par\medskip\noindent\textbf{
\textbf{Observation 6}\textbf{.}
}\par



\emph{For a clause $D$, an assignment $\delta$ of $\text{Vars}(D\circ g)$ falsifies $D\circ g$ if and only if $\overrightarrow{g}(\delta)=\eta_{D}$, i.e., $\overrightarrow{g}(\delta)$ falsifies $D$.}\par



% END-SOURCE-ID Thmtheorem6




\par\medskip\noindent\textbf{
4 Notation}\par




We use the following notation.
For a vector $u$ we use both $u_{i}$ and $u[i]$ to denote the $i$-th entry of $u$, similarly for a matrix $M$ we use $M_{i,j}$ and $M[i;j]$ to denote the entry in the $i$-th row and $j$-th column.
For an ordered set of indices $D$ we denote by $u[D]$ the subvector of $u$ given by $D$, i.e., $u[D]=(u_{i})_{i\in D}$.
For an abbreviation, we use $u[-i]$ to denote the vector $u$ without the $i$-th entry, i.e., $u[-i]=(u_{1},\dots,u_{i-1},u_{i+1},\dots,u_{n})$.\par





\par\medskip\noindent\textbf{
5 Stifling}\par




In this section, we extend the notion of stifling introduced by Chattopadhyay et al. [11].
Let $g:\{0,1\}^{\ell}\to\{0,1\}$ be a Boolean function.
For $i\in[\ell]$ and $b\in\{0,1\}$, we say a partial assignment $\delta\in\{0,1,*\}^{\ell}$ is an $(i,b)$-stifling pattern for $g$ if $\delta_{j}=*$ if and only if $j=i$, and for any $\gamma\in\{0,1\}^{\ell}$ such that $\gamma[-i]=\delta[-i]$, we have $g(\gamma)=b$.
In words, $\delta$ assigns a value to all but the $i$-th bit and when we extend $\delta$ to a full assignment $\gamma$, it holds that $g(\gamma)=b$ no matter how we set the value of the $i$-th bit.\par



% SOURCE-ID Thmtheorem7


\par\medskip\noindent\textbf{
\textbf{Definition 7}\textbf{.}
}\par



\emph{A Boolean function $g:\{0,1\}^{\ell}\to\{0,1\}$ is strongly stifled if there is a collection $P:=\bigl{\{}\delta^{i,b}|i\in[\ell],b\in\{0,1\}\bigr{\}}$ where each $\delta^{i,b}$ is an $(i,b)$-stifling pattern for $g$ and}\par


\[\begin{gathered}
\displaystyle\forall\nobreak\ i\in[\ell],b\in\{0,1\},\text{and }%
\emptyset\neq D\subseteq[\ell]\setminus\{i\} \\
\displaystyle\exists\nobreak\ j\in D\text{ such that }\delta^{j,b}%
\bigl{[}D\setminus\{j\}\bigr{]}=\delta^{i,b}\bigl{[}D\setminus\{j\}\bigr{]}.
\end{gathered}\]


\emph{The collection $P$ is called a converting collection of stifling patterns of $g$.}\par



% END-SOURCE-ID Thmtheorem7



Chattopadhyay et al. [11] defined a stifled function (namely 1-stifled) as a function $g:\{0,1\}^{\ell}\to\{0,1\}$ such that for each $i\in[\ell]$ and $b\in\{0,1\}$ there is an $(i,b)$-stifling pattern for $g$.
In this work, we require not only the existence of the stifling patterns but a stronger property that we can convert the stifling patterns to each other.
More formally, consider an $(i,b)$-stifling pattern $\delta^{i,b}$ from the collection $P$ (from the definition above).
Let an adversary give us a set of coordinates $D\subseteq[\ell]\setminus\{i\}$.
Then, we are able to pick a coordinate $j\in D$ such that the stifling pattern $\delta^{j,b}$ is equal to $\delta^{i,b}$ on all coordinates in $D\setminus\{j\}$.\par




By a simple verification we can show that indexing of two bits $\mathsf{IND}_{1}:\{0,1\}^{3}\to\{0,1\}$ and majority of 3 bits $\mathsf{MAJ}_{3}:\{0,1\}^{3}\to\{0,1\}$ are strongly stifled functions,
where $\mathsf{IND}_{1}(a,d_{0},d_{1})=d_{a}$ and $\mathsf{MAJ}_{3}(x)=1$ if and only if $\bigl{|}\{i\in[3]|x_{i}=1\}\bigr{|}\geq 2$.\par



% SOURCE-ID Thmtheorem8


\par\medskip\noindent\textbf{
\textbf{Observation 8}\textbf{.}
}\par



\emph{The functions $\mathsf{IND}_{1}$ and $\mathsf{MAJ}_{3}$ are strongly stifled.}\par



% END-SOURCE-ID Thmtheorem8



Further, the strongly stifled notion is actually stronger than the original stifled notion, because the inner product function of 2-bit vectors $\mathsf{IP}_{2}:\{0,1\}^{4}\to\{0,1\}$ is stifled [11] but not strongly stifled, where $\mathsf{IP}_{2}(x_{1},x_{2},y_{1},y_{2})=x_{1}x_{2}+y_{1}y_{2}\mod 2$.\par



% SOURCE-ID Thmtheorem9


\par\medskip\noindent\textbf{
\textbf{Observation 9}\textbf{.}
}\par



\emph{The function $\mathsf{IP}_{2}$ is not strongly stifled.}\par



% END-SOURCE-ID Thmtheorem9



For more details, see the appendix.\par





\par\medskip\noindent\textbf{
6 Linear Algebraic Tools}\par




Let $A\subseteq\mathbb{F}_{2}^{m}$ be an affine space over a field $\mathbb{F}_{2}$.
A \emph{constraint representation} of $A$ is a system of linear equations $(M|z)$ where $M\in\mathbb{F}_{2}^{q\times m}$ and $z\in\mathbb{F}_{2}^{q}$ for some $q\leq m$ such that $A=\{y\mid My=z\}$.
The columns of $M$ correspond to the variables of the system $(M|z)$ and rows of $M$ correspond to the constraints.
We say a constraint $i$ \emph{contains} a variable $a$ if $M_{i,a}=1$.
A matrix $M\in\mathbb{F}_{2}^{q\times m}$ is in an \emph{echelon form} if there are $q$ columns $c_{1}<c_{2}<\dots<c_{q}\in[m]$ such that for all $i\in[q]$ it holds that\par


\[\begin{gathered}
M[i;c_{j}]=\begin{cases}1&\text{if $i=j$},\\
0&\text{otherwise}.\end{cases}
\end{gathered}\]


Thus, the submatrix of $M$ induced by the columns $c_{1},\dots,c_{q}$ is the identity matrix $I_{q}\in\mathbb{F}_{2}^{q\times q}$.
The variables corresponding to the columns $c_{1},\dots,c_{q}$ are \emph{dependent} variables of the system $(M|z)$ and the remaining variables are \emph{independent}.
The $c_{i}$-th entry of the $i$-th row of $M$ is called the \emph{pivot of the $i$-th row}.
We say a constraint representation $(M|z)$ is in an \emph{echelon form} if the matrix $M$ is in an echelon form.\par




Let $C\in\mathbb{F}_{2}^{q\times m}$ be a matrix and $t\in\mathbb{F}_{2}^{q}$ be a non-zero vector.
We define a matrix $C^{\prime}=\mathsf{Del}(C,t,i)\in\mathbb{F}_{2}^{q-1\times m}$ where $C^{\prime}$ arises from $C$ by adding the $i$-th row to all rows $j$ such that $t_{j}=1$ and then deleting the row $i$.
Analogously, we define the $\mathsf{Del}$ operation for a constraint representation $(M|z)$ of an affine space $A\subseteq\mathbb{F}_{2}^{m}$, where we treat the vector $z$ as the last column of the matrix $C$.
It turns out that the $\mathsf{Del}$ operation is the only operation needed to get a constraint representation for a super-space as shown in the following theorem that will be the key technical tool for us to process forget nodes while proving our main lifting theorem.\par



% SOURCE-ID Thmtheorem10


\par\medskip\noindent\textbf{
\textbf{Theorem 10}\textbf{.}
}\par



\emph{Let $A_{1}\subseteq A_{2}\subseteq\mathbb{F}_{2}^{m}$ be two affine spaces such that $\dim(A_{2})=\dim(A_{1})+1$.
Let $(M_{1}|z_{1})$ be a constraint representation in the echelon form of $A_{1}$ such that $M_{1}\in\mathbb{F}_{2}^{q\times m}$.
Then, there is a non-zero vector $t\in\mathbb{F}_{2}^{q}$ such that the following is true: for every $i\in[q]$ with $t_{i}=1$, $\mathsf{Del}\bigl{(}(M_{1}|z_{1}),t,i\bigr{)}$ is a constraint representation of $A_{2}$ in echelon form.}\par



% END-SOURCE-ID Thmtheorem10



We call the vector $t$ given by Theorem 10 as \emph{forgetting vector} because it allows us to forget one constraint in the representation of $A_{1}$ to get a representation of $A_{2}$.
Note that the dimension of $t$ equals the number of equations in the system $(M_{1}|z_{1})$.
We say $t$ \emph{contains} a constraint $i$ if $t_{i}=1$.\par




Theorem 10 follows from the following well-known lemma, which we will prove for completeness.
Let $V_{2}\subseteq V_{1}\subseteq\mathbb{F}_{2}^{n}$ be two vector spaces such that $\dim(V_{2})=\dim(V_{1})-1$.\par



% SOURCE-ID Thmtheorem11


\par\medskip\noindent\textbf{
\textbf{Lemma 11}\textbf{.}
}\par



\emph{For any $u,v\in V_{1}\setminus V_{2}$ holds that $u+v\in V_{2}$.}\par



% END-SOURCE-ID Thmtheorem11



\par\medskip\noindent\textbf{Proof.}\par




Since $u$ is in $V_{1}$ but not in $V_{2}$, we have that $\text{Span}(V_{2},u)=V_{1}$ by the dimensions of $V_{1}$ and $V_{2}$.
Since $v\in V_{1}\setminus V_{2}$ as well, we have that $v=w+u$ for an appropriate vector $w\in V_{2}$.
It follows that $u+v=w\in V_{2}$.
$\blacktriangleleft$\par



% SOURCE-ID Thmtheorem12


\par\medskip\noindent\textbf{
\textbf{Corollary 12}\textbf{.}
}\par



\emph{Let $v_{1},\dots,v_{d}$ be a basis of $V_{1}$.
Let $T=\{i\in[d]\mid v_{i}\not\in V_{2}\}$.
Let $j\in T$, and for $i\in[d]\setminus j$ let}\par


\[\begin{gathered}
u_{i}=\begin{cases}v_{i}+v_{j}&\text{if $i\in T$,}\\
v_{i}&\text{otherwise.}\end{cases}
\end{gathered}\]


\emph{Then, $(u_{i})_{i\in[d]\setminus j}$ is a basis of $V_{2}$.}\par



% END-SOURCE-ID Thmtheorem12




\par\medskip\noindent\textbf{Proof.}\par




Vectors $u_{i}$’s are linearly independent as $v_{i}$’s are linearly independent.
By definition and Lemma 11, all $u_{i}$’s are in $V_{2}$.
Since $\dim(V_{2})=d-1$, the vectors $u_{i}$’s has to form a basis of $V_{2}$.
$\blacktriangleleft$\par




Theorem 10 follows from the previous corollary by setting $V_{1}$ and $V_{2}$ to be the constraints spaces of $A_{1}$, and $A_{2}$, respectively (i.e., row spaces of matrices $(M_{1}|z_{1})$ and $(M_{2}|z_{2})$). Further, the forgetting vector of Theorem 10 is the characteristic vector of the set $T$ given by Corollary 12.\par






\par\medskip\noindent\textbf{
7 Simulation}\par




In this section, we prove our lifting theorem.\par




% SOURCE-ID Thmtheorem13


\par\medskip\noindent\textbf{
\textbf{Theorem 13}\textbf{ }(Theorem 1 stated for F(P)DT)\textbf{.}
}\par



\emph{Let $R\subseteq\{0,1\}^{n}\times O$ be a relation and $\mathcal{T}$ be a canonical FPDT computing $R\circ g$ where $g:\{0,1\}^{\ell}\to\{0,1\}$ is a strongly stifled gadget.
Then, there is an FDT $\mathcal{T}^{\prime}$ computing $R$ such that $d_{\mathsf{q}}(\mathcal{T}^{\prime})\leq\log|\mathcal{T}|_{\mathsf{q}}$ and $w(\mathcal{T}^{\prime})\leq w(\mathcal{T})$.}\par



% END-SOURCE-ID Thmtheorem13




\par\medskip\noindent\textbf{Algorithm}\par




We prove Theorem 13 by simulation.
On an input $x\in\{0,1\}^{n}$, the constructed FDT $\mathcal{T}^{\prime}$ simulates given FPDT $\mathcal{T}$ on an input $y\in\{0,1\}^{m}$ (for $m:=n\ell$) with $x=\overrightarrow{g}(y)$ by traversing a path from the root $r$ of $\mathcal{T}$ to a leaf.
The main loop of the simulation is quite simple.
We start in the root $r$ of $\mathcal{T}$ and in each iteration, we process the current node $v$ of $\mathcal{T}$ and pick a new node.
Sometimes during the processing of a node of $\mathcal{T}$, we query or forget a bit of $x$.
When we reach a leaf $s$ of $\mathcal{T}$ we just output the value of $s$.
The main loop is summarized in Algorithm 1.\par




Let $v$ be a current node of $\mathcal{T}$ we have just encountered.
We maintain a constraint representation in echelon form $(M_{v}|z_{v})$ of the affine space $A_{v}$.
We store queried (and not forgotten) bits of $x$ in a partial assignment $\rho_{v}\in\{0,1,*\}^{n}$.
Let $C(\rho_{v})\subseteq\{0,1\}^{n}$ be a set of all possible extensions of $\rho_{v}$ and $w(\rho_{v})$ be a number of fixed bits of $\rho_{v}$.
Thus, $C(\rho_{v})$ is a cube and $w(\rho_{v})$ is its width.
Our goal is that any $\bar{x}\in C(\rho_{w})$ is represented in $A_{v}$ (i.e., there is $y\in A_{v}$ such that $\overrightarrow{g}(y)=\bar{x}$), and that $w(\rho_{v})$ is at most the co-dimension of $A_{v}$.\par




An input $y$ of $\mathcal{T}$ is divided into $n$ blocks $B_{1},\dots,B_{n}\subseteq[m]$, each of size $\ell$, and each such block corresponds to exactly one entry of $x$.
Formally, $B_{i}=\{(i-1)\ell+1,\dots,i\ell\}$.
During the simulation of $\mathcal{T}$, we divide the blocks into two groups – free and fixed.
Fixed blocks correspond to the entries of $x$ that were queried and were not forgotten – i.e., the entries fixed by $\rho_{v}$.
The other blocks are free.\par




For each fixed bit $i\in[n]$ of $\rho_{v}$, we have a unique constraint $j_{i}$ of $(M_{v}|z_{v})$ such that the pivot of the constraint $j_{i}$ is in the block $B_{i}$.
The constraint $j_{i}$ is called \emph{the primary constraint of $B_{i}$}.
The dependent variables of primary constraints are called \emph{marked variables}.
We will keep an invariant that each marked variable is in a different block, i.e., each fixed block contains a unique marked variable.
The other (non-marked) variables of the fixed blocks are called \emph{stifling variables}.
We will assign values to the stifling variables according to appropriate stifling patterns, in such a way that a fixed block $B_{i}$ is always mapped to $x_{i}$, independently on the value of the marked variable of $B_{i}$.
Further, we will assign a value to the marked variable in a block $B_{i}$ in such a way that the assignment will satisfy the corresponding primary constraint.
Since each marked variable is in a unique primary constraint, we will show that such an assignment of all marked variables always exists.\par




The constraints that are not primary for any block are called \emph{secondary}.
It will turn out that all secondary constraints contain variables only from fixed blocks.
All variables of free blocks are called \emph{free}.
Thus, the matrix $M_{v}$ has the following form (after rearranging columns):\par







Let $P:=\bigl{\{}\delta^{i,b}\mid i\in[\ell],b\in\{0,1\}\bigr{\}}$ be a converting collection of stifling patterns of $g$, given by the assumption.
Let $\alpha_{v}\in\{0,1,*\}^{m}$ be the following partial assignment:\par


\[\begin{gathered}
\alpha_{v}[B_{i}]=\begin{cases}\delta^{j,x_{i}}&\text{if $B_{i}$ is a fixed %
block and $j$ is the index of the marked variable of $B_{i}$}\\
*^{\ell}&\text{if $B_{i}$ is a free block}\end{cases}
\end{gathered}\]


Note that the partial assignments $\alpha_{v}$ assigns values exactly to the stifling variables of $M_{v}$.
We will keep an invariant that if we set all stifling variables according to $\alpha_{v}$ all secondary constraints of $(M_{v}|z_{v})$ will be satisfied.
This will help us to ensure that each $\bar{x}\in C(\rho_{v})$ is represented in $A_{v}$.\par




At the beginning of our simulation, we are at the root $r$ of $\mathcal{T}$.
Since $A_{r}=\{0,1\}^{m}$, the matrix $M_{r}$ is an empty matrix and all variables are free because we have not queried any entry of $x$ yet.
Also, the patterns $\rho_{r}$ and $\alpha_{r}$ do not contain any fixed bit, i.e. they are equal to $*^{n}$, or $*^{m}$, respectively.\par




\textbf{Algorithm 1}  Simulation.


\textbf{Input:} $x\in\{0,1\}^{n}$  $\triangleright$ Input for FDT $\mathcal{T}^{\prime}$\par





FPDT $\mathcal{T}$ with the root $r$ computing $R\circ g$\par



 



\textbf{Initialization:}\par





 



\textbf{Simulation:}\par








During the simulation, we will maintain the following invariants.
This will help us to ensure that each $\bar{x}\in C(\rho_{v})$ is represented in $A_{v}$.\par


\textbf{ Invariant 1:}



The system of equations $(M_{v}|z_{v})$ is a constraint representation of $A_{v}$ in the echelon form.\par



\textbf{ Invariant 3:}



All variables of all free blocks are independent.\par



\textbf{ Invariant 5:}



For each fixed bit $i\in[n]$ of $\rho_{v}$, there is a unique constraint $j_{i}$ of $(M_{v}|z_{v})$ such that the pivot (i.e., the marked variable) of the constraint $j_{i}$ is in the block $B_{i}$.\par



\textbf{ Invariant 7:}



The partial assignment $\alpha_{v}$ assigns values to all stifling variables and any extension of $\alpha_{v}$ to a full assignment satisfies all secondary constraints.\par







Note that Invariant \textbf{ Invariant 7:} implies that secondary constraints of $(M_{v}|z_{v})$ contain only variables of fixed blocks.
We will show these invariants hold for any node $v$ of $\mathcal{T}$ at the moment, when $v$ is checked whether $v$ is a leaf – i.e., at Line 4 of Algorithm 1.
Clearly, the invariants hold for the root $r$ of $\mathcal{T}$.
Now, we describe how we process a node $v$ (depending on whether $v$ is a query node or forget node).
We suppose the invariants hold for $v$.
During the processing, we pick an appropriate child $u$ of $v$ and make $u$ the new current node.
Subsequently, we argue why the invariants hold for $u$.\par




We remark that the query node processing is a careful adaptation of the node processing given by Chattopadhyay et al. [11].
All new machinery (strongly stifled function and obtaining a constraint representation of a super-space by Theorem 10) is used only for the processing of forget nodes.\par





\par\medskip\noindent\textbf{Query Nodes}\par




When $v$ is a query node, then $v$ introduces a new parity query $f_{v}$, and if $\langle f_{v},y\rangle=0$ the computation of $\mathcal{T}$ proceeds to the left child $u_{0}$ of $v$, otherwise to the right child $u_{1}$.
Our goal is to pick an appropriate child $u$ of $v$ and create the system $(M_{u}|z_{u})$ representing $A_{u}$ satisfying all our requirements.
Let us start with a system $\bigl{(}M^{\prime}|z(b)\bigr{)}$, where\par


\[\begin{gathered}
M^{\prime}=\begin{pmatrix}M_{v}\\
f_{v}\end{pmatrix},z(b)=\begin{pmatrix}z_{v}\\
b\end{pmatrix},
\end{gathered}\]


with $b$ being a parameter equal to 0 or 1.
We fix the value of $b$ when we pick the appropriate child $u$ of $v$ as the new node.
Surely, the system $\bigl{(}M^{\prime}|z(b)\bigr{)}$ represents the space $A_{u_{b}}$, however, it might not satisfy our requirements (for example the matrix $M^{\prime}$ might not be in the echelon form).
Note that the matrix $M^{\prime}$ does not depend on the value of $b$.
We do another pivoting step of the Gaussian elimination to get the system $\bigl{(}M^{\prime}|z(b)\bigr{)}$ into the echelon form, i.e.,\par


\begin{enumerate}


\item[] 
1. 


We zero out all coefficients in $f_{v}$ corresponding the dependent variables in $(M_{v}|z_{v})$, to get a new constraint $(f^{\prime}|b^{\prime})$, where $b^{\prime}$ is a function of $b$. We call the new constraint $(f^{\prime}|b^{\prime})$ the \emph{reduced form} of the constraint $(f_{v}|b)$.\par





\item[] 
2. 


We pick one of the remaining variables $a$ contained in $f^{\prime}$ as a new dependent variable, we pick an appropriate child $u$ of $v$ and we set the value of $b$ (and $b^{\prime}$), accordingly.\par





\item[] 
3. 


We zero out all coefficients corresponding to $a$ in all original constraints from the system $(M_{v}|z_{v})$ to get the new system $(M_{u}|z_{u})$.\par





\end{enumerate}


It is clear the new system $(M_{u}|z_{u})$ is a constraint representation of $A_{u}$ (i.e., Invariant \textbf{ Invariant 1:} will hold for $u$).
The crucial part is to pick a new dependent variable $a$ in Step 2 of the executed Gaussian elimination.
Note that the reduced constraint $(f^{\prime}|b^{\prime})$ does not contain any marked variable as all marked variables are dependant and, thus, they are zeroed out from $(f|b)$ in Step 1 of the executed pivoting step.
There are two cases to consider as follows.\par




\par\medskip\noindent\textbf{Case 1.}\par




The new constraint $(f^{\prime}|b^{\prime})$ contains only variables of the fixed blocks.
Then, the new constraint becomes a secondary constraint, and the new dependent variable $a$ can be any variable of $f^{\prime}$.
Since the constraint $(f^{\prime}|b^{\prime})$ contains only variables of fixed blocks (but no marked variable), we can assign a value to all variables of $(f^{\prime}|b^{\prime})$ given by $\alpha_{v}$.
Thus, there is $\bar{b}\in\{0,1\}$ such that for any extension $y$ of $\alpha_{v}$, it holds that $\langle f^{\prime},y\rangle=\bar{b}$.
Then, we pick the appropriate child $u$ of $v$, that gives us the right value of $b$ (and $b^{\prime}$) such that the new constraint $(f^{\prime}|b^{\prime})$ holds for any extension of $\alpha_{v}$. This ensures that Invariant \textbf{ Invariant 7:} holds for $u$.\par




We did not query any new bit of $x$ in this case.
It follows that the partial assignment $\rho_{v}$ and the set of fixed and free blocks are not changed.
The set of primary constraints is unchanged as well.
Further, the set of pivots of $(M_{v}|z_{v})$ is not changed by the pivoting step of the Gaussian elimination.
Thus, Invariant \textbf{ Invariant 5:} holds for $u$.
Invariant \textbf{ Invariant 3:} holds because the constraint $(f^{\prime}|b^{\prime})$ does not contain any free variable of $(M_{v}|z_{v})$ and thus the new dependent variable $a$ can not be from a free block.\par





\par\medskip\noindent\textbf{Case 2.}\par




The new constraint $(f^{\prime}|b^{\prime})$ contains at least one variable $a$ of a free block $B_{i}$.
In this case, we can pick the new vertex $u$ as an arbitrary child of $v$.
Let $\mathcal{T}_{w}$ be a subtree of $\mathcal{T}$ rooted in a node $w$ of $\mathcal{T}$.
We compare the query size of subtrees $\mathcal{T}_{u_{0}}$ and $\mathcal{T}_{u_{1}}$ and we pick $u$ to be the root of the subtree with the smaller query size, i.e., $|\mathcal{T}_{u}|_{\mathsf{q}}\leq|\mathcal{T}_{w}|_{\mathsf{q}}$, where $w$ is the other children of $v$.\par




We query $x_{i}$ and update the partial assignment $\rho_{v}$ by the value of $x_{i}$ to get $\rho_{u}$.
The block $B_{i}$ becomes a fixed block.
The new constraint $(f^{\prime}|b^{\prime})$ becomes the primary constraint of $B_{i}$ and the variable $a\in B_{i}$ becomes the pivot of $(f^{\prime}|b^{\prime})$, i.e., $a$ becomes a marked variable.
Since the set of pivots of $(M_{v}|z_{v})$ is not changed, Invariant \textbf{ Invariant 5:} holds for $u$.
Since the only new dependant variable is $a\in B_{i}$, Invariant \textbf{ Invariant 3:} holds as well.\par




The partial assignment $\alpha_{u}$ differs only at the block $B_{i}$ from $\alpha_{v}$ ($\alpha_{v}[B_{i}]=*^{\ell}$, and $\alpha_{u}[B_{i}]=\delta^{j,x_{i}}$).
Since the block $B_{i}$ was free in $(M_{v}|z_{v})$, no secondary constraint of $(M_{v}|z_{v})$ contains any variable of the block $B_{i}$.
Thus, no secondary constraints of $(M_{v}|z_{v})$ were changed by the pivoting step in this case.
The new constraint $(f^{\prime}|b^{\prime})$ is primary.
Thus, no secondary constraint of $(M_{u}|z_{u})$ contains any variable of the block $B_{i}$ as well.
Therefore, any extension of $\alpha_{u}$ still satisfies all secondary constraints of $(M_{u}|z_{u})$ and Invariant \textbf{ Invariant 7:} holds for $\alpha_{u}$.\par




See Algorithm 2 for a summary of the query node processing.\par



\textbf{Algorithm 2}  \textbf{Process Query Node} ($v$: the current (query) node).





\par\medskip\noindent\textbf{Forget Nodes}\par




In the case when $v$ is a forget node, the node $v$ has the only child $u$ and $\dim(A_{u})=\dim(A_{v})+1$.
We have the constraint representation $(M_{v}|z_{v})$ of $A_{v}$ maintained by our simulation for $M_{v}\in\mathbb{F}_{2}^{c\times m}$ and $z_{v}\in\mathbb{F}_{2}^{c}$ (where $c$ is the co-dimension of $A_{v}$).
For processing the forget node, we introduce a classification of stifling variables.
The variables of fixed blocks that are contained in the secondary constraints are called dangerous\footnote{Our notion of dangerous \emph{variables} is not to be confused with the notion of dangerous \emph{constraints} used by Efremenko et al. [13].}.
Note that the marked variables can not be dangerous.
The remaining variables of fixed blocks (i.e., non-marked and non-dangerous) are called safe.
Thus, with this new classification, the matrix $M_{v}$ has the following form:\par







Let $t\in\mathbb{F}_{2}^{c}$ be a forgetting vector given by an application of Theorem 10 to spaces $A_{v}$ and $A_{u}$.
The new system $(M_{u}|z_{u})$ is obtained after applying $\mathsf{Del}\bigl{(}(M_{v}|z_{v}),t,i)$ for a right choice of $i$ (the function $\mathsf{Del}$ is defined in Section 6).
By Theorem 10, the system $(M_{u}|z_{u})$ is a constraint representation of $A_{u}$ in the echelon form, i.e. Invariant \textbf{ Invariant 1:} holds for $u$.
Let $p$ be the number of primary constraints in $(M_{v}|z_{v})$, i.e. wlog, the constraints $1,\dots,p$ are primary and the constraints $p+1,\dots,c$ are secondary.
We consider two cases.\par




\par\medskip\noindent\textbf{Case 1.}\par




There is an $i\in\{p+1,\dots,c\}$ such that $t_{i}=1$.
Then, fix one such $i$ and take a system $(M_{u}|z_{u})=\mathsf{Del}\bigl{(}(M_{v}|z_{v}),t,i)$.
We do not query or forget any bit of $x$, thus $\rho_{u}=\rho_{v}$ and $\alpha_{u}=\alpha_{v}$.
To create $(M_{u}|z_{u})$, we only added the secondary constraint $i$ to some rows of $(M_{v}|z_{v})$ and then we deleted it.
Thus, the set of variables which appear in secondary constraints cannot grow in size and, therefore, secondary constraints are still satisfied by the assignment $\alpha_{u}$.
Therefore, Invariant \textbf{ Invariant 7:} holds for $\alpha_{u}$.\par




The set of primary constraints is not changed.
The $\mathsf{Del}$ operation does not change the set of marked variables as the secondary constraint $i$ does not contain any pivot of the primary constraints.
Thus, Invariants \textbf{ Invariant 5:} holds for $u$.
The set of fixed blocks does not change and there is no new dependant variable. Thus, Invariant \textbf{ Invariant 3:} holds as well.\par





\par\medskip\noindent\textbf{Case 2.}\par




For all $i\in\{p+1,\dots,c\}$, it holds that $t_{i}=0$.
Then, we fix some $i\in\{1,\dots,p\}$ such that $t_{i}=1$.
Note that such an $i$ exists as $t$ is a non-zero vector.
Again, let $(M_{u}|z_{u})=\mathsf{Del}\bigl{(}(M_{v}|z_{v}),t,i)\bigr{)}$.
Since $t$ has only zeroes at the coordinates corresponding to the secondary constraints, the secondary constraints are not changed by the $\mathsf{Del}$ operation.
As the constraint $i$ is deleted and it was a primary one, one marked variable $a$ (the pivot of the constraint $i$) becomes independent and safe.
Let $B_{j}$ be the block containing the variable $a$, i.e., the constraint $i$ of $(M_{v}|z_{v})$ is the primary constraint for $B_{j}$.
We consider two sub-cases.\par





\par\medskip\noindent\textbf{Sub-case 2.1.}\par




The other variables of $B_{j}$ are safe in $(M_{u}|z_{u})$ as well, i.e., they are not in any secondary constraint.
Thus, the whole block $B_{j}$ contains only independent and safe variables of $(M_{u}|z_{u})$.
We forget the bit $x_{j}$ and make the block $B_{j}$ free.
The set of other primary constraints (different from $i$) may change their form, but their pivots are not changed. Hence, Invariant \textbf{ Invariant 5:} holds for $u$.
There is no new dependant variable. Thus Invariant \textbf{ Invariant 3:} holds as well.\par




We get the partial assignment $\rho_{u}$ by simply setting the variable $x_{j}$ free.
The partial assignment $\alpha_{u}$ differs from $\alpha_{v}$ only at the block $B_{j}$ ($\alpha_{u}[B_{j}]=*^{\ell}$, and $\alpha_{v}[B_{j}]=\delta^{j,x_{j}}$).
Further, the secondary constraints of $(M_{u}|z_{u})$ do not contain any variable of the block $B_{j}$ by the assumption.
Thus, Invariant \textbf{ Invariant 7:} holds for $\alpha_{u}$.\par





\par\medskip\noindent\textbf{Sub-case 2.2.}\par




There is a dangerous variable in the block $B_{j}$, i.e., there is a secondary constraint of $(M_{u}|z_{u})$ that contains a variable of $B_{j}$.
In this case, we use the strong stifling property of $g$.
Let $D\subseteq[\ell]$ be the set of indices of all dangerous variables of $(M_{u}|z_{u})$ in $B_{j}$.
Let $j_{1}$ be the index of the variable $a$ in $B_{j}$ (i.e., the previously marked variable in $B_{j}$).
Note that $j_{1}\not\in D$ because the variable $a$ is safe.
Thus by definition of strongly stifling, there is a $j_{2}\in D$ such that $\delta^{j_{2},x_{j}}\bigl{[}D\setminus\{j_{2}\}\bigr{]}=\delta^{j_{1},x_{j}}%
\bigl{[}D\setminus\{j_{2}\}\bigr{]}$ (note that $\alpha_{v}[B_{j}]=\delta^{j_{1},x_{j}}$).
Let $a^{\prime}$ be the $j_{2}$-th variable in the block $B_{j}$ and $k$ be a secondary constraint that contains $a^{\prime}$ (such constraint exists by the assumption).\par




We run again the pivoting step for $a^{\prime}$, i.e., we zero out all coefficients corresponding to $a^{\prime}$ in all other constraints of $(M_{u}|z_{u})$ by adding the constraints $k$ to all other constraints containing $a^{\prime}$.
We denote the final system of constraints as $(M^{\prime}_{u}|z^{\prime}_{u})$.
Note that $(M^{\prime}_{u}|z^{\prime}_{u})$ is still a constraint representation of $A_{u}$ as it arises from $(M_{u}|z_{u})$ only by row operations.\par




The constraint $k$ is now the only constraint containing the variable $a^{\prime}$ and $a^{\prime}$ becomes a dependent variable.
Thus, we make the constraint $k$ a primary constraint for $B_{j}$ and we mark the variable $a^{\prime}$.
The primary constraint for $B_{j}$ was changed from the constraint $i$ of $(M_{v}|z_{v})$ to the constraint $k$ of $(M^{\prime}_{u}|z^{\prime}_{u})$ and the marked variable in the block $B_{j}$ was changed from $a$ to $a^{\prime}$.
The set of other primary constraints and their pivots were not changed.
Thus, Invariant \textbf{ Invariant 5:} holds for $u$.\par




We do not change the assignment $\rho_{v}$, thus the sets of free and fixed blocks are the same.
The only change in the set of dependent variables was done in the block $B_{j}$ (that remains a fixed block), thus Invariant \textbf{ Invariant 3:} holds for $u$.\par




The secondary constraints of $(M_{v}|z_{v})$ were not changed by the $\mathsf{Del}\bigl{(}(M_{v}|z_{v}),t,i)$ executed at the beginning of this case (as $t_{i^{\prime}}=0$ for all secondary constraints $i^{\prime}$).
Since $k$ is a secondary constraint of $(M_{u}|z_{u})$, the secondary constraints of $(M^{\prime}_{u}|z^{\prime}_{u})$ contains only variables of fixed blocks.
However, we change the marked variable in the block $B_{j}$.
Thus, the partial assignment $\alpha_{u}$ differs from $\alpha_{v}$ at the block $B_{j}$ ($\alpha_{v}=\delta^{j_{1},x_{i}}$, and $\alpha_{u}=\delta^{j_{2},x_{i}}$, where $j_{1}$ and $j_{2}$ are indices of $a$ and $a^{\prime}$ in the block $B_{j}$).
We need to be sure that $\alpha_{u}$ still gives a solution to the secondary constraints of $(M^{\prime}_{u}|z^{\prime}_{u})$.
Note that the secondary constraints of $(M^{\prime}_{u}|z^{\prime}_{u})$ might still contain variables from the block $B_{j}$.\par




By pivoting $a^{\prime}$ and making the constraint $k$ primary, the variable $a^{\prime}$ is not in any secondary constraint of $(M^{\prime}_{u}|z^{\prime}_{u})$.
Since $k$ was a secondary constraint of $(M_{u}|z_{u})$, it can not happen that a safe variable in $(M_{u}|z_{u})$ would become a dangerous one in $(M^{\prime}_{u}|z^{\prime}_{u})$ (i.e., by the pivoting of $a^{\prime}$).
In other words, the set of variables of the secondary constraints of $(M^{\prime}_{u}|z^{\prime}_{u})$ is a subset of the set of variables of the secondary constraints of $(M_{u}|z_{u})$.
Thus, the set $D\setminus\{j_{2}\}$ still contains all dangerous variables of $B_{j}$ in $(M^{\prime}_{u}|z^{\prime}_{u})$.
Since $\alpha_{v}\bigl{[}D\setminus\{j_{2}\}\bigr{]}=\alpha_{u}\bigl{[}D\setminus\{j_%
{2}\}\bigr{]}$ by the assumption, any extension $\alpha_{u}$ satisfy all secondary constraints of $(M^{\prime}_{u}|z^{\prime}_{u})$ and Invariant \textbf{ Invariant 7:} holds for $\alpha_{u}$.\par




A summary of the forget node processing is in Algorithm 3.\par



\textbf{Algorithm 3}  \textbf{Process Forget Node} ($v$: the current (forget) node).





\par\medskip\noindent\textbf{Proof of Theorem 13
}\par




Theorem 13 follows from the following lemma.\par



% SOURCE-ID Thmtheorem14


\par\medskip\noindent\textbf{
\textbf{Lemma 14}\textbf{.}
}\par



\emph{Suppose the simulation is at Line 4 of Algorithm 1, i.e., it checks whether the current node $v$ is a leaf.
Then,}\par


\begin{enumerate}


\item[] 
1. 


$w\bigl{(}C(\rho_{v})\bigr{)}\leq\text{co-dim}(A_{v})$\emph{.}\par





\item[] 
2. 


\emph{For any }$\bar{x}\in C(\rho^{v})$\emph{, there is }$y\in A_{v}$\emph{ such that }$\overrightarrow{g}(y)=x$\emph{.}\par





\end{enumerate}



% END-SOURCE-ID Thmtheorem14



\par\medskip\noindent\textbf{Proof of Item 1.}\par




By Invariant \textbf{ Invariant 1:}, the co-dimension of $A_{v}$ is exactly the number of equations in the system $(M_{v}|z_{v})$.
By Invariant \textbf{ Invariant 5:}, the number of fixed bits by $\rho_{v}$ is exactly the number of primary constraints in $(M_{v}|z_{v})$.
Thus, $w\bigl{(}C(\rho_{v})\bigr{)}\leq\text{co-dim}(A_{v})$.
$\blacktriangleleft$\par





\par\medskip\noindent\textbf{Proof of Item 2.}\par




Let $\bar{x}\in C(\rho_{v})$.
We will find a solution $y$ to the system $(M_{v}|z_{v})$ such that $\overrightarrow{g}(y)=\bar{x}$.
Thus, by Invariant \textbf{ Invariant 1:}, $y\in A_{v}$.\par




First, we set variables of free blocks.
Let $B_{i}$ be a free block.
Thus, by Invariant \textbf{ Invariant 3:}, all variables of $B_{i}$ are independent.
We set the variables of $B_{i}$ in a way such that the block $B_{i}$ is mapped to $\bar{x}_{i}$ by the gadget $g$.\par




Now, we set the values of the stifling variables according to $\alpha_{v}$.
By Invariant \textbf{ Invariant 7:}, all secondary constraints are satisfied by any extension of $\alpha_{v}$.
Recall that for a fixed block $B_{i}$, $\alpha_{v}[B_{i}]=\delta^{j,x_{i}}$ where $j$ is the index of the marked variable of $B_{i}$ and $\bar{x}_{i}=\rho_{v}[i]$.
Since $\delta^{j,x_{i}}$ is a $(j,x_{i})$-stifling pattern, it holds that the block $B_{i}$ will be always mapped to $\bar{x}_{i}$ by $g$, no matter how we set the marked variables.
Thus, the constructed solution $y$ will be mapped onto $\bar{x}$.
By Invariant \textbf{ Invariant 5:}, each primary constraint contains a unique marked variable.
Thus, we can set a value to each marked variable $a$ in such a way the primary constraint containing $a$ is satisfied.
$\blacktriangleleft$\par





\par\medskip\noindent\textbf{Proof of Theorem 13.}\par




By Item 1 of Lemma 14, the width of a cube $C(\rho_{v})$ in a time of checking whether a vertex $v$ is a leaf is at most co-dimension of $A_{v}$.
Thus, the width of the constructed FDT $\mathcal{T}^{\prime}$ is at most the width of $\mathcal{T}$.\par




Now, we bound the query depth of $\mathcal{T}^{\prime}$.
Consider a root-leaf path $P$ of $\mathcal{T}^{\prime}$ and let $d$ be the number of queries made on $P$.
Note that any time we query a bit of $x$ (Line 9 of Algorithm 2) we also pick a subtree with a smaller query size (Line 8 of Algorithm 2).
Thus, by each query of $\mathcal{T}^{\prime}$ we halve the query size of $\mathcal{T}$.
Thus, $2^{d}\leq|\mathcal{T}|_{\mathsf{q}}$.\par




It remains to prove the constructed FDT $\mathcal{T}^{\prime}$ is correct.
Let $s$ be a leaf of $\mathcal{T}$ that is reached during the simulation and $o\in O$ is the output of $s$.
Since $\mathcal{T}$ computes $R\circ g$, it holds that for all $y\in A_{s}$ we have $(y,o)\in R\circ g$.
Note that the processing phase (Lines 5-8 of Algorithm 1) is not executed for any leaf.
Thus, the assertion of Lemma 14 holds for the leaf $s$ even at the time of output – Line 9 of Algorithm 1.
Therefore at the end of the simulation, for any $\bar{x}\in C(\rho_{s})$ there is $y\in A_{s}$ such that $\overrightarrow{g}(y)=\bar{x}$.
Since $(y,o)\in R\circ g$, it holds that $(\bar{x},o)\in R$ and the constructed FDT $\mathcal{T}^{\prime}$ indeed outputs a correct answer.
$\blacktriangleleft$\par







\par\medskip\noindent\textbf{
8 Application}\par




Razborov [24] showed the following trade-off between width and size of tree-like resolution.\par



% SOURCE-ID Thmtheorem15


\par\medskip\noindent\textbf{
\textbf{Theorem 15}\textbf{ }(Theorem 3.1, Razborov [24])\textbf{.}
}\par



\emph{Let $k=k(n)\geq 4$ be any parameter and let $\varepsilon>0$ be an arbitrary constant.
Then, there exists a $k$-CNF contradiction $\tau^{\prime}$ over $n$ variables such that there is a resolution refutation for $\tau^{\prime}$ with width at most $O(k)$, but for any tree-like resolution refutation $\Pi$ for $\tau^{\prime}$ with $w(\Pi)\leq n^{1-\varepsilon}/k$, we have the bound $|\Pi|\geq\exp(n^{\Omega(k)})$.}\par



% END-SOURCE-ID Thmtheorem15



By our simulation, given by Theorem 13, we can lift the trade-off (given by the previous theorem) to tree-like $\text{Res}(\oplus)$ and prove Theorem 3.\par




% SOURCE-ID Thmtheorem3


\par\medskip\noindent\textbf{
\textbf{Theorem 3}\textbf{.}
[Restated, see original statement.]}\par



\emph{Let $k=k(n)\geq 12$ be any parameter and let $\varepsilon>0$ be an arbitrary constant.
Then, there exists a $k$-CNF contradiction $\tau$ over $n$ variables such that there is a resolution refutation for $\tau$ with width at most $O(k)$, but for every tree-like $\text{Res}(\oplus)$ refutation $\Pi$ for $\tau$ with $w(\Pi)\leq n^{1-\varepsilon}/k$, we have the bound $|\Pi|\geq\exp(n^{\Omega(k)})$.}\par



% END-SOURCE-ID Thmtheorem3




\par\medskip\noindent\textbf{Proof.}\par




Let $g:\{0,1\}^{3}\to\{0,1\}$ be a strongly stifled gadget – such functions exist as observed in Section 5.
Let $k^{\prime}:=\lfloor k/3\rfloor$, and $\tau^{\prime}$ be a $k^{\prime}$-CNF contradiction given by Theorem 15.
We set $\tau:=\tau^{\prime}\circ g$ that is a $k$-CNF contradiction.
Since there is a resolution refutation for $\tau^{\prime}$ with width at most $O(k^{\prime})$, then there is a resolution refutation for $\tau$ with width at most $O(k)$.\par




Now, let $\Pi$ be a tree-like $\text{Res}(\oplus)$ refutation for $\tau$ with $w(\Pi)\leq n^{1-\varepsilon}/k$.
By Observation 4, let $\mathcal{T}$ be a canonical FPDT corresponding to $\Pi$ that computes $\textit{Search}(\tau)$.
Thus, we have $w(\mathcal{T})\leq w(\Pi)+1$ and $|\mathcal{T}|_{q}\leq|\Pi|$.
We change $\mathcal{T}$ to compute $\textit{Search}(\tau^{\prime})\circ g$.
Let $s$ be a leaf of $\mathcal{T}$ outputting a clause $D$ of $\tau^{\prime}\circ g$.
The clause $D$ has to appear in a set of clauses $D^{\prime}\circ g$ for a clause $D^{\prime}$ of $\tau^{\prime}$.
We change the output of $s$ to be the clause $D^{\prime}$ instead of $D$.
By Observation 6, the tree $\mathcal{T}$ now computes $\textit{Search}(\tau^{\prime})\circ g$.\par




By Theorem 13, there is FDT $\mathcal{T}^{\prime}$ computing $\textit{Search}(\tau^{\prime})$ with $d_{\mathsf{q}}(\mathcal{T}^{\prime})\leq\log|\mathcal{T}|_{\mathsf{q}}$ and $w(\mathcal{T}^{\prime})\leq w(\mathcal{T})$.
Let $\Pi^{\prime}$ be the resolution refutation for $\tau^{\prime}$ corresponding to the succinct form of $\mathcal{T}^{\prime}$.
Thus, $w(\Pi^{\prime})=w(\mathcal{T}^{\prime})$ and $|\Pi^{\prime}|\leq 3\cdot 2^{d_{\mathsf{q}}(\mathcal{T}^{\prime})}$ (by Observation 5).
Since\par


\[\begin{gathered}
w(\mathcal{T}^{\prime})\leq w(\mathcal{T})\leq n^{1-\varepsilon}/k+1\leq n^{1-%
\varepsilon}/k^{\prime},
\end{gathered}\]


we have that $|\Pi^{\prime}|\geq\exp(n^{\Omega(k^{\prime})})$ by Theorem 15.
The last inequality in (1) holds if $k\leq 2n^{1-\varepsilon}$, which holds as we suppose that $1\leq w(\Pi)\leq n^{1-\varepsilon}/k$.
Putting everything together, we have\par


\[\begin{gathered}
|\Pi|\geq|\mathcal{T}|_{\mathsf{q}}\geq 2^{d_{\mathsf{q}}(\mathcal{T}^{\prime}%
)}\geq\frac{1}{3}\cdot|\Pi^{\prime}|\geq\exp(n^{\Omega(k^{\prime})})=\exp(n^{%
\Omega(k)}).\
\end{gathered}\]


$\blacktriangleleft$\par






\par\medskip\noindent\textbf{References}\par



\begin{itemize}


\item[] 
[1]

Yaroslav Alekseev and Dmitry Itsykson.


Lifting to regular resolution over parities via games.


\emph{Electron. Colloquium Comput. Complex.}, TR24-128, 2024.


URL: https://eccc.weizmann.ac.il/report/2024/128.





\item[] 
[2]

Sepehr Assadi, Gillat Kol, and Zhijun Zhang.


Rounds vs communication tradeoffs for maximal independent sets.


In \emph{63rd IEEE Annual Symposium on Foundations of Computer
Science, FOCS 2022, Denver, CO, USA, October 31 - November 3, 2022}, pages
1193–1204. IEEE, 2022.


doi:10.1109/FOCS54457.2022.00115.





\item[] 
[3]

Paul Beame, Christopher Beck, and Russell Impagliazzo.


Time-space tradeoffs in resolution: superpolynomial lower bounds for
superlinear space.


In Howard J. Karloff and Toniann Pitassi, editors, \emph{Proceedings
of the 44th Symposium on Theory of Computing Conference, STOC 2012, New
York, NY, USA, May 19 - 22, 2012}, pages 213–232. ACM, 2012.


doi:10.1145/2213977.2213999.





\item[] 
[4]

Paul Beame and Sajin Koroth.


On disperser/lifting properties of the index and inner-product
functions.


In Yael Tauman Kalai, editor, \emph{14th Innovations in Theoretical
Computer Science Conference, ITCS 2023, January 10-13, 2023, MIT,
Cambridge, Massachusetts, USA}, volume 251 of \emph{LIPIcs}, pages
14:1–14:17. Schloss Dagstuhl – Leibniz-Zentrum für Informatik, 2023.


doi:10.4230/LIPICS.ITCS.2023.14.





\item[] 
[5]

Chris Beck, Jakob Nordström, and Bangsheng Tang.


Some trade-off results for polynomial calculus: extended abstract.


In Dan Boneh, Tim Roughgarden, and Joan Feigenbaum, editors, \emph{Symposium on Theory of Computing Conference, STOC’13, Palo Alto, CA, USA,
June 1-4, 2013}, pages 813–822. ACM, 2013.


doi:10.1145/2488608.2488711.





\item[] 
[6]

Eli Ben-Sasson.


Size space tradeoffs for resolution.


In John H. Reif, editor, \emph{Proceedings on 34th Annual ACM
Symposium on Theory of Computing, May 19-21, 2002, Montréal,
Québec, Canada}, pages 457–464. ACM, 2002.


doi:10.1145/509907.509975.





\item[] 
[7]

Christoph Berkholz, Moritz Lichter, and Harry Vinall-Smeeth.


Supercritical size-width tree-like resolution trade-offs for graph
isomorphism, 2024.


doi:10.48550/arXiv.2407.17947.





\item[] 
[8]

Christoph Berkholz and Jakob Nordström.


Supercritical space-width trade-offs for resolution.


\emph{SIAM J. Comput.}, 49(1):98–118, 2020.


doi:10.1137/16M1109072.





\item[] 
[9]

Sreejata Kishor Bhattacharya, Arkadev Chattopadhyay, and Pavel Dvořák.


Exponential Separation Between Powers of Regular and General
Resolution over Parities.


In Rahul Santhanam, editor, \emph{39th Computational Complexity
Conference (CCC 2024)}, volume 300 of \emph{Leibniz International Proceedings
in Informatics (LIPIcs)}, pages 23:1–23:32, Dagstuhl, Germany, 2024. Schloss
Dagstuhl – Leibniz-Zentrum für Informatik.


doi:10.4230/LIPIcs.CCC.2024.23.





\item[] 
[10]

Mark Braverman and Rotem Oshman.


A rounds vs. communication tradeoff for multi-party set disjointness.


In Chris Umans, editor, \emph{58th IEEE Annual Symposium on
Foundations of Computer Science, FOCS 2017, Berkeley, CA, USA, October
15-17, 2017}, pages 144–155. IEEE Computer Society, 2017.


doi:10.1109/FOCS.2017.22.





\item[] 
[11]

Arkadev Chattopadhyay, Nikhil S. Mande, Swagato Sanyal, and Suhail Sherif.


Lifting to parity decision trees via stifling.


In Yael Tauman Kalai, editor, \emph{14th Innovations in Theoretical
Computer Science Conference, ITCS 2023, January 10-13, 2023, MIT,
Cambridge, Massachusetts, USA}, volume 251 of \emph{LIPIcs}, pages
33:1–33:20. Schloss Dagstuhl – Leibniz-Zentrum für Informatik, 2023.


doi:10.4230/LIPICS.ITCS.2023.33.





\item[] 
[12]

Susanna F. de Rezende, Noah Fleming, Duri Andrea Janett, Jakob Nordström, and
Shuo Pang.


Truly supercritical trade-offs for resolution, cutting planes,
monotone circuits, and weisfeiler-leman, 2024.


doi:10.48550/arXiv.2411.14267.





\item[] 
[13]

Klim Efremenko, Michal Garlík, and Dmitry Itsykson.


Lower bounds for regular resolution over parities.


In Bojan Mohar, Igor Shinkar, and Ryan O’Donnell, editors, \emph{Proceedings of the 56th Annual ACM Symposium on Theory of Computing, STOC
2024, Vancouver, BC, Canada, June 24-28, 2024}, pages 640–651. ACM, 2024.


doi:10.1145/3618260.3649652.





\item[] 
[14]

Noah Fleming, Toniann Pitassi, and Robert Robere.


Extremely Deep Proofs.


In Mark Braverman, editor, \emph{13th Innovations in Theoretical
Computer Science Conference (ITCS 2022)}, volume 215 of \emph{Leibniz
International Proceedings in Informatics (LIPIcs)}, pages 70:1–70:23,
Dagstuhl, Germany, 2022. Schloss Dagstuhl – Leibniz-Zentrum für
Informatik.


doi:10.4230/LIPIcs.ITCS.2022.70.





\item[] 
[15]

Lance Fortnow.


Time-space tradeoffs for satisfiability.


\emph{J. Comput. Syst. Sci.}, 60(2):337–353, 2000.


doi:10.1006/jcss.1999.1671.





\item[] 
[16]

Lance Fortnow, Richard J. Lipton, Dieter van Melkebeek, and Anastasios Viglas.


Time-space lower bounds for satisfiability.


\emph{J. ACM}, 52(6):835–865, 2005.


doi:10.1145/1101821.1101822.





\item[] 
[17]

Ankit Garg, Mika Göös, Pritish Kamath, and Dmitry Sokolov.


Monotone circuit lower bounds from resolution.


\emph{Theory Comput.}, 16:1–30, 2020.


Preliminary version in STOC 2018.


doi:10.4086/TOC.2020.V016A013.





\item[] 
[18]

Mika Göös, Gilbert Maystre, Kilian Risse, and Dmitry Sokolov.


Supercritical tradeoffs for monotone circuits, 2024.


doi:10.48550/arXiv.2411.14268.





\item[] 
[19]

Dmitry Itsykson and Dmitry Sokolov.


Lower bounds for splittings by linear combinations.


In Erzsébet Csuhaj-Varjú, Martin Dietzfelbinger, and
Zoltán Ésik, editors, \emph{Mathematical Foundations of Computer
Science 2014 - 39th International Symposium, MFCS 2014, Budapest, Hungary,
August 25-29, 2014. Proceedings, Part II}, volume 8635 of \emph{Lecture
Notes in Computer Science}, pages 372–383. Springer, 2014.


doi:10.1007/978-3-662-44465-8\_32.





\item[] 
[20]

Dmitry Itsykson and Dmitry Sokolov.


Resolution over linear equations modulo two.


\emph{Ann. Pure Appl. Log.}, 171(1), 2020.


doi:10.1016/j.apal.2019.102722.





\item[] 
[21]

Noam Nisan and Avi Wigderson.


Rounds in communication complexity revisited.


\emph{SIAM J. Comput.}, 22(1):211–219, 1993.


doi:10.1137/0222016.





\item[] 
[22]

Pavel Pudlák.


On extracting computations from propositional proofs (a survey).


In Kamal Lodaya and Meena Mahajan, editors, \emph{IARCS Annual
Conference on Foundations of Software Technology and Theoretical Computer
Science, FSTTCS 2010, December 15-18, 2010, Chennai, India}, volume 8 of
\emph{LIPIcs}, pages 30–41. Schloss Dagstuhl – Leibniz-Zentrum für
Informatik, 2010.


doi:10.4230/LIPIcs.FSTTCS.2010.30.





\item[] 
[23]

Alexander A. Razborov.


Unprovability of lower bounds on circuit size in certain fragments of
bounded-arithmetic.


\emph{Izvestiya. Math.}, 59(1):205–227, 1995.





\item[] 
[24]

Alexander A. Razborov.


A new kind of tradeoffs in propositional proof complexity.


\emph{J. ACM}, 63(2):16:1–16:14, 2016.


doi:10.1145/2858790.





\item[] 
[25]

Dmitry Sokolov.


Dag-like communication and its applications.


In Pascal Weil, editor, \emph{Computer Science – Theory and
Applications – 12th International Computer Science Symposium in Russia,
CSR 2017, Kazan, Russia, June 8-12, 2017, Proceedings}, volume 10304 of
\emph{Lecture Notes in Computer Science}, pages 294–307. Springer, 2017.


doi:10.1007/978-3-319-58747-9\_26.





\item[] 
[26]

R. Ryan Williams.


Time-space tradeoffs for counting NP solutions modulo integers.


\emph{Comput. Complex.}, 17(2):179–219, 2008.


doi:10.1007/s00037-008-0248-y.





\end{itemize}




\par\medskip\noindent\textbf{
Appendix A Appendix}\par




In this section, we show that the functions $\mathsf{IND}_{1}$ and $\mathsf{MAJ}_{3}$ are strongly stifled and $\mathsf{IP}_{2}$ is not strongly stifled.\par




% SOURCE-ID Thmtheorem8


\par\medskip\noindent\textbf{
\textbf{Observation 8}\textbf{.}
[Restated, see original statement.]}\par



\emph{The functions $\mathsf{IND}_{1}$ and $\mathsf{MAJ}_{3}$ are strongly stifled.}\par



% END-SOURCE-ID Thmtheorem8




\par\medskip\noindent\textbf{Proof.}\par




We present collections of $(i,b)$-stifling patterns $P(\mathsf{IND}_{1})$ and $P(\mathsf{MAJ}_{3})$ for $\mathsf{IND}_{1}$ and $\mathsf{MAJ}_{3}$, respectively.
It is straight-forward to verify that these collections are converting collections of stifling patterns for $\mathsf{IND}_{1}$ and $\mathsf{MAJ}_{3}$.\par






Table 1: $P(\mathsf{IND}_{1})$.




0
1




1
$(*,0,0)$
$(*,1,1)$


2
$(1,*,0)$
$(1,*,1)$


3
$(0,0,*)$
$(0,1,*)$







Table 2: $P(\mathsf{MAJ}_{3})$.




0
1




1
$(*,0,0)$
$(*,1,1)$


2
$(0,*,0)$
$(1,*,1)$


3
$(0,0,*)$
$(1,1,*)$









$\blacktriangleleft$\par




% SOURCE-ID Thmtheorem9


\par\medskip\noindent\textbf{
\textbf{Observation 9}\textbf{.}
[Restated, see original statement.]}\par



\emph{The function $\mathsf{IP}_{2}$ is not strongly stifled.}\par



% END-SOURCE-ID Thmtheorem9





\par\medskip\noindent\textbf{Proof.}\par




The only $(1,1)$-stifling pattern for $\mathsf{IP}_{2}:\{0,1\}^{4}\to\{0,1\}$ is $\delta^{1}:=(*,1,0,1)$.
Similarly, the only $(2,1)$- and $(4,1)$-stifling patterns for $\mathsf{IP}_{2}$ are $\delta^{2}:=(1,*,1,0)$, and $\delta^{4}:=(1,0,1,*)$, respectively.
Now, let $D=\{2,4\}$.
There is no $j\in D$ such that $\delta^{1}\bigl{[}D\setminus\{j\}\bigr{]}=\delta^{j}\bigl{[}D\setminus\{j\}%
\bigr{]}$ as required to be a strongly stifled function.
$\blacktriangleleft$\par




\par




\end{document}
